\metatitle{Non-commutative symmetric functions}
\metadescription{An introduction to non-commutative symmetric functions, including non-commutative Schur functions and a non-commutative Murnaghan--Nakayama rule.}
\metakeywords{Non-commutative symmetric functions,non-commutative Schur functions,Murnaghan--Nakayama rule} 

\section[nonComFunctions]{Non-commutative symmetric functions}

For an introduction to non-commutative symmetric functions, 
we refer to \cite{GelfandKrobLascouxLeclercRetakhThibon1995}.

The non-commutative symmetric functions are dual to the quasisymmetric functions.
These are not to be confused with \hyperref[schurInNonCom]{symmetric functions in non-commutative variables}.
One convenient model is
\[
  \mathrm{NSym}=\setZ\langle \completeH_1,\completeH_2,\dotsc\rangle,
\]
the free associative algebra generated by noncommutative complete homogeneous
functions.  For a composition $\alpha$, write
$\completeH_\alpha\coloneqq
\completeH_{\alpha_1}\completeH_{\alpha_2}\dotsm \completeH_{\alpha_\ell}$.
The pairing with quasisymmetric functions is normalized by
$\langle \completeH_\alpha,\qmonom_\beta\rangle=\delta_{\alpha,\beta}$.
There is a forgetful morphism $\chi:\mathrm{NSym}\to\spaceSym$ sending
$\completeH_n$ to the ordinary complete homogeneous symmetric function.

\begin{polydata}{nonComHomogeneous}
  Name   & Noncommutative complete homogeneous functions \\
  Space    & NSym \\
  Basis    & Yes \\
  Rating   & 3 \\
  Bib      & GelfandKrobLascouxLeclercRetakhThibon1995 \\
  Symbol   & $\completeH_\alpha$ \\
  Year     & 1995 \\
  Category & Elementary \\
  PositiveIn & nonComPowerPsi | BallantineDaughertyHicksMason2020 | proof=dual to the $\qPsi$ monomial expansion \\
  PositiveIn & nonComPowerPhi | BallantineDaughertyHicksMason2020 | proof=dual to the $\qPhi$ monomial expansion \\
  NSymQSymDual & qmonom | GelfandKrobLascouxLeclercRetakhThibon1995 | pairing=canonical NSym--QSym pairing \\
\end{polydata}

\begin{polydata}{nonComPowerPsi}
  Name   & Noncommutative power sums (Psi) \\
  Space    & NSym \\
  Basis    & Yes \\
  Rating   & 2 \\
  Bib      & GelfandKrobLascouxLeclercRetakhThibon1995,BallantineDaughertyHicksMason2020 \\
  Symbol   & $\Psi_\alpha$ \\
  Year     & 1995 \\
  Category & Elementary \\
  NSymQSymDual & qPsi | BallantineDaughertyHicksMason2020 | pairing=canonical NSym--QSym pairing \\
\end{polydata}

\begin{polydata}{nonComPowerPhi}
  Name   & Noncommutative power sums (Phi) \\
  Space    & NSym \\
  Basis    & Yes \\
  Rating   & 2 \\
  Bib      & GelfandKrobLascouxLeclercRetakhThibon1995,BallantineDaughertyHicksMason2020 \\
  Symbol   & $\Phi_\alpha$ \\
  Year     & 1995 \\
  Category & Elementary \\
  NSymQSymDual & qPhi | BallantineDaughertyHicksMason2020 | pairing=canonical NSym--QSym pairing \\
\end{polydata}

The two noncommutative power-sum bases above are dual to the two
\hyperref[qPsi]{quasisymmetric power-sum bases} $\qPsi$ and $\qPhi$.
Since $\qPsi_\alpha$ and $\qPhi_\alpha$ have positive expansions in the
\hyperref[qmonom]{monomial quasisymmetric basis}, duality gives positive
expansions of $\completeH_\alpha$ in the corresponding noncommutative
power-sum bases \cite{BallantineDaughertyHicksMason2020}.

\name{Angela Hicks} and \name{Robert McCloskey} give a concrete
realization of $\mathrm{NSym}$ in noncommuting variables and derive the
standard noncommutative elementary, complete, ribbon, and power-sum bases from
that model \cite{HicksMcCloskey2024x}.  Their account also gives
combinatorial change-of-basis formulas using minimal products of transition
matrices and brick tabloid statistics, paralleling the classical symmetric
function picture.

\name{John M. Campbell} introduces a lift of Stanley's chromatic symmetric
function to $\mathrm{NSym}$ \cite{Campbell2024x}.  For an unlabeled directed
graph $D$, the construction gives an element $X_D\in \mathrm{NSym}$ whose
commutative image is the ordinary chromatic symmetric function of the
underlying undirected graph.  The lift is obtained from Stanley's power-sum
expansion using the $\Psi$-basis of $\mathrm{NSym}$, and it leads to
digraph-indexed generating sets for $\mathrm{NSym}$.
\name[J.-C. Novelli]{Jean-Christophe Novelli} and
\name[J.-Y. Thibon]{Jean-Yves Thibon} study chromatic quasisymmetric
functions inside the algebra of quasisymmetric functions in noncommuting
variables \cite{NovelliThibon2025x}.  Their work also proposes
noncommutative Macdonald polynomials compatible with a Haglund--Wilson type
formula and relates the construction to Yang--Baxter elements in Hecke
algebras.

\name{Edward E. Allen} and \name{Sarah Katherine Mason} give a combinatorial
formula for the expansion of immaculate noncommutative symmetric functions in
the complete homogeneous noncommutative basis \cite{AllenMason2025}.  Their
model uses GBPR diagrams and tunnel hooks, extending the role played by special
rim hooks in the classical inverse Kostka formula.

\name{Diego Arcis}, \name{Camilo Gonz{\'a}lez}, and
\name{Sebasti{\'a}n M{\'a}rquez} study the Hopf algebra of noncommutative
symmetric functions in superspace \cite{ArcisGonzalezMarquez2024}.  They
construct primitive elements, extend the elementary and power-sum families to
superspace, introduce noncommutative ribbon Schur functions in superspace, and
realize the Hopf algebra using trees.
\name{Franz Lehner}, \name{Jean-Christophe Novelli}, and
\name{Jean-Yves Thibon} study combinatorial Hopf algebras in
noncommutative probability \cite{LehnerNovelliThibon2020x}.
\name{Li Guo}, \name{Jean-Yves Thibon}, and \name{Houyi Yu} study Hopf
algebras of signed permutations, weak quasisymmetric functions, and
Malvenuto--Reutenauer type structures \cite{GuoThibonYu2020}.

\section[nonComSchur]{Non-commutative Schur functions}

\begin{polydata}{nonComSchur}
  Name   & Non-commutative Schur functions \\
  Space    & NSym \\
  Basis    & Yes \\
  Rating   & 1 \\
  Bib      & GelfandKrobLascouxLeclercRetakhThibon1995 \\
  Keywords & non-commutative \\
  Year     & 1995 \\
  Category & Schur \\
\end{polydata}

There are several Schur-like bases in $\mathrm{NSym}$.  The
\defin{noncommutative Schur functions} of
\name{Christine Bessenrodt}, \name{Kurt Luoto}, and
\name{Stephanie van Willigenburg} are the basis dual to the
\hyperref[qSchur]{quasisymmetric Schur functions}
\cite{BessenrodtLuotoWilligenburg2011}.  Equivalently, they give a
composition-indexed lift of ordinary Schur functions under the forgetful map
$\chi:\mathrm{NSym}\to\spaceSym$.  The older ribbon basis of
\cite{GelfandKrobLascouxLeclercRetakhThibon1995} is another fundamental
Schur-like basis of $\mathrm{NSym}$.

\section[immaculateSchur]{Immaculate Schur functions}

\begin{polydata}{immaculateSchur}
  Name     & Immaculate Schur functions \\
  Space    & NSym \\
  Basis    & Yes \\
  Rating   & 1 \\
  Bib      & BergBergeronSaliolaSerranoZabrocki2014 \\
  Year     & 2014 \\
  Symbol   & $\mathfrak{S}_\alpha$ \\
  Category & Schur \\
\end{polydata}

The \defin{immaculate Schur functions} form a composition-indexed basis of
$\mathrm{NSym}$ introduced by \name{Chris Berg}, \name{Nantel Bergeron},
\name{Franco Saliola}, \name{Luis Serrano}, and \name{Mike Zabrocki}
\cite{BergBergeronSaliolaSerranoZabrocki2014}.  They are constructed using
Bernstein-like creation operators, and the forgetful map
$\chi:\mathrm{NSym}\to\spaceSym$ sends $\mathfrak{S}_\lambda$ to the ordinary
\hyperref[schurS]{Schur function} $\schurS_\lambda$ when $\lambda$ is a
partition.

The immaculate basis has a positive right Pieri rule and a Jacobi--Trudi type
formula.  Its dual basis is the \hyperref[diSchur]{dual immaculate Schur
functions} in \hyperref[quasiSymmetricFunctions]{QSym}.

\section[nonComMNrule]{Murnaghan--Nakayama rule}

\name{Vasu Tewari} proves a noncommutative analogue of the classical
\hyperref[murnaghanNakayamaRule]{Murnaghan--Nakayama rule} for the
noncommutative Schur functions of \name{Christine Bessenrodt},
\name{Kurt Luoto}, and \name{Stephanie van Willigenburg} \cite{Tewari2016}.
The rule expands a product of a noncommutative power sum and a
noncommutative Schur function into noncommutative Schur functions, with signs
controlled by noncommutative border strips.
